\exercisesection{\S~3}

\begin{enumerate}
\def\labelenumi{\arabic{enumi})}
\item
  设 g 是李代数，\(g'\) 是 g 的嘉当子代数。则命题 3 的条件满足。但 \(g'\) 的元素即使在 \(g'\) 中正则，也未必在 g 中正则。
\item
  设 g 是 n 维实李代数，U（相应地 H）是 g 的正则元集合（相应地嘉当子代数集合），\(\text{Int}(\mathfrak{g})\) 是 g 的内自同构群（第三章 §6 第 2 节 定义 2）。
\end{enumerate}

\begin{enumerate}
\def\labelenumi{\alph{enumi})}
\item
  证明：若 \(x\) 与 \(y\) 属于 \(\mathrm{U}\) 的同一连通分支，则 \(\mathfrak{g}^0(x)\) 与 \(\mathfrak{g}^0(y)\) 在 \(\operatorname{Int}(\mathfrak{g})\) 下共轭。
\item
  证明 U 的连通分支个数有限，且该个数被一个只依赖于 n 的常数 \(c(n)\) 所界（应用附录 II 习题 2）。
\item
  推出 \(\operatorname{Int}(\mathfrak{g})\) 在 \(\mathrm{H}\) 上的轨道数 \(\leq c(n)\) 。
\end{enumerate}

\begin{enumerate}
\def\labelenumi{\arabic{enumi})}
\setcounter{enumi}{2}
\tightlist
\item
  设 \(\mathfrak{g}\) 是实李代数，\(\mathfrak{r}\) 是 \(\mathfrak{g}\) 的根基，\(\mathfrak{h}\) 与 \(\mathfrak{h}'\) 是 \(\mathfrak{g}\) 的嘉当子代数，\(\varphi\) 是从 \(\mathfrak{g}\) 到 \(\mathfrak{g} / \mathfrak{r}\) 的典则同态。下列条件等价：
\end{enumerate}

\begin{enumerate}
\def\labelenumi{(\roman{enumi})}
\item
  \(\mathfrak{h}\) 与 \(\mathfrak{h}'\) 在 \(\operatorname{Int}(\mathfrak{g})\) 下共轭；
\item
  \(\varphi(\mathfrak{h})\) 与 \(\varphi(\mathfrak{h}')\) 在 \(\operatorname{Int}(\mathfrak{g}/\mathfrak{r})\) 下共轭。（模仿命题 5 的证明。）
\end{enumerate}

\begin{enumerate}
\def\labelenumi{\arabic{enumi})}
\setcounter{enumi}{3}
\item
  设 \(\mathfrak{g} = \mathfrak{sl}(2, \mathbf{R})\) ，\(x = \begin{pmatrix} 1 & 0 \\ 0 & -1 \end{pmatrix}\) ，\(y = \begin{pmatrix} 0 & -1 \\ 1 & 0 \end{pmatrix}\) 。证明 \(\mathbf{R}x\) 与 \(\mathbf{R}y\) 是 \(\mathfrak{g}\) 的嘉当子代数，且它们在 \(\operatorname{Aut}(\mathfrak{g})\) 下不共轭。
\item
  \begin{enumerate}
  \def\labelenumii{\alph{enumii})}
  \tightlist
  \item
    证明存在李代数 \(\mathfrak{g}\) （ \(k\) 上的），具有基 \((x,y,z,t)\) ，使得
  \end{enumerate}
\end{enumerate}

\[
[ x, y ] = z, [ x, t ] = t, [ y, t ] = 0, [ \mathfrak {g}, z ] = 0.
\]

证明 g 可解，且 \(k = kx + ky + kz\) 是 g 的子代数。

\begin{enumerate}
\def\labelenumi{\alph{enumi})}
\setcounter{enumi}{1}
\item
  证明 g 的初等自同构恰是形如 \(1 + \lambda \operatorname{ad}_{\mathfrak{g}} y + \mu \operatorname{ad}_{\mathfrak{g}} t\) 的映射，其中 \(\lambda, \mu \in k\) 。
\item
  证明 \(1 + ad_{k}x\) 是 k 的初等自同构，它不能开拓为 g 的初等自同构。
\end{enumerate}

\(d\) ) 设 \(\mathfrak{s}\) 是李代数 \(\mathfrak{a}\) 的半单子代数。证明 \(\mathfrak{s}\) 的每个初等自同构都可开拓为 \(\mathfrak{a}\) 的初等自同构。

\begin{enumerate}
\def\labelenumi{\arabic{enumi})}
\setcounter{enumi}{5}
\item
  约化李代数 \(\mathfrak{g}\) 的每个元素都含于一个维数为 \(\mathrm{rk}(\mathfrak{g})\) 的交换子代数中。
\item
  设 \(\mathfrak{g}\) 是李代数，\(\mathfrak{g}'\) 是 \(\mathfrak{g}\) 的子代数且在 \(\mathfrak{g}\) 中约化。设 \(\mathfrak{a}\) 是 \(\mathfrak{g}'\) 的嘉当子代数。证明存在 \(\mathfrak{g}\) 的一个嘉当子代数包含 \(\mathfrak{a}\) （利用 §2 命题 10）。推出 \(\mathrm{rk}(\mathfrak{g}') \leq \mathrm{rk}(\mathfrak{g})\) ，且等号成立当且仅当 \(\mathfrak{g}'\) 具有命题 3 的性质 (i)、(ii)、(iii)。
\item
  设 \(\mathfrak{g}\) 是李代数，\(\mathfrak{a}\) 是 \(\mathfrak{g}\) 的理想，\(\mathfrak{h}\) 是 \(\mathfrak{g}\) 的嘉当子代数，\(\mathcal{C}_{\infty}\mathfrak{g}\) 是 \(\mathfrak{g}\) 的升中心列之并。证明 \(\mathfrak{a} \subset \mathfrak{h}\) 蕴涵 \(\mathfrak{a} \subset \mathcal{C}_{\infty}\mathfrak{g}\) （换言之，\(\mathcal{C}_{\infty}\mathfrak{g}\) 是 \(\mathfrak{g}\) 中含于 \(\mathfrak{h}\) 的最大理想）。（化归到 \(k\) 代数闭的情形，并注意 \(\mathfrak{a}\) 在 \(\mathfrak{g}\) 的每个初等自同构下稳定；于是关系 \(\mathfrak{a} \subset \mathfrak{h}\) 蕴涵 \(\mathfrak{a}\) 含于 \(\mathfrak{g}\) 的每个嘉当子代数；借助 §2 习题 15 得出结论。）
\end{enumerate}

¶ 9) 设 g 是李代数，h 是 g 的嘉当子代数，x 是 h 的元素。设 \(g = h \oplus g^{+}\) 是 g 关于 h 的 Fitting 分解（§1 第 1 节）。

\begin{enumerate}
\def\labelenumi{\alph{enumi})}
\item
  设 \(\mathfrak{n}\) 是最大的半单 \(\mathfrak{h}\) -子模，它含于 \(\mathfrak{g}^0 (x)\cap \mathfrak{g}^+\) 。证明 \(\mathfrak{n} = 0\) 当且仅当 \(\mathfrak{g}^0 (x) = \mathfrak{h}\) ，即 \(x\) 在 \(\mathfrak{g}\) 中正则。
\item
  证明 \(\mathfrak{h} \oplus \mathfrak{n}\) 是 \(\mathfrak{g}\) 的子代数。若 \(\mathfrak{h}'\) 是 \(\mathfrak{h}\) 与 \(\mathfrak{n}\) 的交换子之交，证明 \(\mathfrak{h}'\) 是 \(\mathfrak{h} \oplus \mathfrak{n}\) 的理想，它包含 \(\mathcal{D}\mathfrak{h}\) 与 \(x\) 。由此得出（习题 8）\(\mathfrak{h}' \subset \mathcal{C}_{\infty}(\mathfrak{h} \oplus \mathfrak{n})\) ，从而 \(x \in \mathcal{C}_{\infty}(\mathfrak{h} \oplus \mathfrak{n})\) ，且 \(x\) 属于 \(\mathfrak{h} \oplus \mathfrak{n}\) 的每个嘉当子代数。
\item
  若 \(n \neq 0\) ，则 \(h \oplus n\) 不幂零且有无穷多个嘉当子代数（§2 习题 10）。由此得出 x 属于 g 的无穷多个嘉当子代数。
\end{enumerate}

\(d\) ) \(\mathfrak{g}\) 的元素是正则元，当且仅当它属于唯一的嘉当子代数。

¶ 10) 设 g 是李代数，r 是其根基，n 是其最大幂零理想，a 是其一个 Levi 子代数。

\begin{enumerate}
\def\labelenumi{\alph{enumi})}
\item
  置 \(\mathfrak{g}' = \mathfrak{n} + \mathcal{D}\mathfrak{g}\) 。证明 \(\mathfrak{g}' = \mathfrak{n} \oplus \mathfrak{a}\) 。（利用 \([\mathfrak{g}, \mathfrak{r}]\) 含于 \(\mathfrak{n}\) 这一事实。）若 \(\mathfrak{g} \neq 0\) ，则 \(\mathfrak{g}' \neq 0\) 。
\item
  假设 k 代数闭。设 \((\mathrm{V}_{i})_{i\in\mathrm{I}}\) 是 g-模 g（取伴随表示）的 Jordan-Hölder 序列的商。若 \(x\in r\) ，证明 \(x_{V_{i}}\) 是位似，且 \(x_{V_{i}}=0\) 对一切 i 成立当且仅当 x 属于 n。推出元素 \(y\in g\) 属于 \(g'\) 当且仅当 \(\operatorname{Tr}(y_{\mathrm{V}_{i}})=0\) 对一切 \(i\in I\) 成立。
\item
  用 N 表示 g 中由使 ad x 幂零的元素 x 生成的向量子空间。证明 N 是 g 的子代数（利用 N 在 \(\operatorname{Aut}_{e}(\mathfrak{g})\) 下稳定这一事实）。利用 b) 证明 \(N \subset g'\) 。
\end{enumerate}

\(d\) ) 设 \(\mathfrak{h}\) 是 \(\mathfrak{g}\) 的嘉当子代数。假设存在子集 \(R\) （\(\mathfrak{h}^*\) 中的），使得

\[
\mathfrak {g} = \mathfrak {h} \oplus \bigoplus_ {\alpha \in R} \mathfrak {g} ^ {\alpha} (\mathfrak {h}),
\]

这一假定特别在 k 代数闭时满足。证明 N 此时包含诸 \(\mathfrak{g}^{\alpha}(\mathfrak{h})\) 、Dh 与 n；推出 N 包含 \(g'\) ，故 \(N = g'\) 。

\begin{enumerate}
\def\labelenumi{\alph{enumi})}
\setcounter{enumi}{4}
\tightlist
\item
  若 \(k\) 代数闭且 \(\mathfrak{g} \neq 0\) ，则 \(\mathfrak{g}\) 含有元素 \(x \neq 0\) 使得 \(\operatorname{ad} x\) 幂零。（事实上，此时有 \(\mathfrak{g}' \neq 0\) 。）
\end{enumerate}

¶ 11) 设 g 是李代数，r 是其根基。

\begin{enumerate}
\def\labelenumi{\alph{enumi})}
\item
  设 \(\mathfrak{s}\) 是 \(\mathfrak{g}\) 的 Levi 子代数，\(\mathfrak{k}\) 是 \(\mathfrak{s}\) 的嘉当子代数。证明 \(\mathfrak{k}\) 含于一个嘉当子代数 \(\mathfrak{h}\) （\(\mathfrak{g}\) 中的），后者是 \(\mathfrak{k}\) 与 \(\mathfrak{r}\) 的一个子代数之和。（利用 §2 定理 2、命题 10 及定理 1 的推论 2。）
\item
  设 \(\mathfrak{h}'\) 是 \(\mathfrak{g}\) 的嘉当子代数。证明存在 Levi 子代数 \(\mathfrak{s}'\) （\(\mathfrak{g}\) 的），使得 \(\mathfrak{h}'\) 是 \(\mathfrak{s}'\) 的一个嘉当子代数与 \(\mathfrak{r}\) 的一个子代数之和。（子代数 \(\mathfrak{s}, \mathfrak{k}, \mathfrak{h}\) （\(a\) ) 中的）可以选得使 \(\mathfrak{h} + \mathfrak{r} = \mathfrak{h}' + \mathfrak{r}\) 。置 \(\mathfrak{a} = \mathfrak{h} + \mathfrak{r}\) ，它是可解的。由定理 3，存在 \(x \in \mathcal{C}^{\infty}(\mathfrak{a})\) 使得 \(e^{\mathrm{ad}_{\mathfrak{a}}x}\mathfrak{h} = \mathfrak{h}'\) 。于是 \(e^{\mathrm{ad}_{\mathfrak{g}}x}\) 是 \(\mathfrak{g}\) 的特殊自同构，它把 \(\mathfrak{s}\) 变为所需的 Levi 子代数。）
\item
  设 \(\mathfrak{s}\) 是 \(\mathfrak{g}\) 的 Levi 子代数。设 \(\mathfrak{h}\) 是 \(\mathfrak{g}\) 的嘉当子代数，它是嘉当子代数 \(\mathfrak{k}\) （\(\mathfrak{s}\) 的）与子代数 \(\mathfrak{l}\) （\(\mathfrak{r}\) 的）之和。设 \(\mathfrak{c}\) 是 \(\mathfrak{k}\) 在 \(\mathfrak{r}\) 中的交换子。证明 \(\mathfrak{l}\) 是 \(\mathfrak{c}\) 的嘉当子代数。（对 \(x \in \mathfrak{k}\) ，\(\operatorname{ad}_{\mathfrak{g}} x\) 半单，而 \(\operatorname{ad}_{\mathfrak{h}} x\) 幂零，故 \([\mathfrak{k}, \mathfrak{h}] = 0\) 。若 \(y \in \mathfrak{c}\) 使得 \([y, \mathfrak{l}] \subset \mathfrak{l}\) ，则 \([y, \mathfrak{h}] \subset \mathfrak{h}\) ，故 \(y \in \mathfrak{h} \cap \mathfrak{r} = \mathfrak{l}\) 。）
\item
  设 \(\mathfrak{s}\) 是 \(\mathfrak{g}\) 的 Levi 子代数，\(\mathfrak{k}\) 是 \(\mathfrak{s}\) 的嘉当子代数，\(\mathfrak{c}\) 是 \(\mathfrak{k}\) 在 \(\mathfrak{r}\) 中的交换子，\(\mathfrak{l}\) 是 \(\mathfrak{c}\) 的嘉当子代数。则 \(\mathfrak{h} = \mathfrak{k} + \mathfrak{l}\) 是 \(\mathfrak{g}\) 的嘉当子代数。（设 \(x = y + z\) （ \(y \in \mathfrak{s}, z \in \mathfrak{r}\) ）是正规化子 \(\mathfrak{u}\) （\(\mathfrak{h}\) 在 \(\mathfrak{g}\) 中的）的元素。证明 \([y, \mathfrak{k}] \subset \mathfrak{k}\) ，故 \(y \in \mathfrak{k}\) 且 \(z \in \mathfrak{u}\) 。再证明 \([z, \mathfrak{k}] \subset \mathfrak{h} \cap \mathfrak{r} \subset \mathfrak{c}\) ，故 \([\mathfrak{k}, [\mathfrak{k}, z]] = 0\) ，从而 \([\mathfrak{k}, z] = 0\) 且 \(z \in \mathfrak{c}\) 。最后，\([z, \mathfrak{l}] \subset \mathfrak{l}\) ，故 \(z \in \mathfrak{l}\) 且 \(x \in \mathfrak{h}\) 。）
\item
  设 \(\mathfrak{s},\mathfrak{k},\mathfrak{c}\) 如 \(d\) ) 中所示，\(\mathfrak{q}=[\mathfrak{g},\mathfrak{r}]\) 是 \(\mathfrak{g}\) 的幂零根基。设 \(x\in\mathfrak{q}\) ，\(u\) 是特殊自同构 \(e^{\mathrm{ad}x}\) 。若 \(u(\mathfrak{k})\subset\mathfrak{k}+\mathfrak{c}\) ，则 \(x\in\mathfrak{c}\) 。（考察 \(\rho\) ，即 \(\mathfrak{s}\) 在 \(\mathfrak{q}\) 上的伴随表示，设 \(\mathfrak{q}_{i}\) 是 \(\mathcal{C}^{i+1}\mathfrak{q}\) 在 \(\mathcal{C}^{i}\mathfrak{q}\) 中的在 \(\rho\) 下稳定的补；设 \(\rho_{i}\) 是 \(\rho\) 的由 \(\mathfrak{q}_{i}\) 定义的子表示。设 \(\sigma_{i}=\rho_{i}|\mathfrak{k}\) 。设 \(\mathfrak{q}_{i}^{\prime}\) 是 \(\mathfrak{k}\) 在 \(\mathfrak{q}_{i}\) 中的交换子，\(\mathfrak{q}_{i}^{\prime\prime}\) 是 \(\mathfrak{q}_{i}^{\prime}\) 在 \(\mathfrak{q}_{i}\) 中的在 \(\sigma_{i}\) 下稳定的补。设 \(x=x_{1}^{\prime}+x_{1}^{\prime\prime}+\cdots+x_{n}^{\prime}+x_{n}^{\prime\prime}\) ，其中 \(x_{i}^{\prime}\in\mathfrak{q}_{i}^{\prime},x_{i}^{\prime\prime}\in\mathfrak{q}_{i}^{\prime\prime}\) 。用反证法，设诸 \(x_{i}^{\prime\prime}\) 不全为零，例如 \(x_{1}^{\prime\prime}=\cdots=x_{p-1}^{\prime\prime}=0\) ，\(x_{p}^{\prime\prime}\neq0\) 。若 \(h\in\mathfrak{k}\) ，则 \(u(h)=h+[x_{p}^{\prime\prime},h]+y\) ，其中 \(y\in\mathcal{C}^{p+1}\mathfrak{q}\) 。由于 \(u(h)\in\mathfrak{k}+\mathfrak{c}\) ，这给出 \([x_{p}^{\prime\prime},h]+y\in\mathfrak{q}_{p}^{\prime}+\mathfrak{q}_{p+1}^{\prime}+\cdots+\mathfrak{q}_{n}^{\prime}\) ，从而 \([h,x_{p}^{\prime\prime}]\in\mathfrak{q}_{p}^{\prime}\) ，故 \([h,x_{p}^{\prime\prime}]=0\) 。于是 \(x_{p}^{\prime\prime}=0\) ，矛盾。）
\end{enumerate}

\(f\) ) 设 \(\mathfrak{h}\) 是 \(\mathfrak{g}\) 的嘉当子代数。则 \(\mathfrak{h}\) 可以唯一地表示为 \(\mathfrak{h} \cap \mathfrak{r}\) 与 \(\mathfrak{g}\) 的某个 Levi 子代数的嘉当子代数之和。（关于唯一性，利用 \(e\) ) 及第一章 §6 第 8 节定理 5。）所述 Levi 子代数一般并不唯一。

\(g\) ) 设 \(\mathfrak{h}\) 是 \(\mathfrak{g}\) 的嘉当子代数，并设 \(\mathfrak{t} = \mathfrak{g}^0 (\mathfrak{h} \cap \mathfrak{r})\) 。则 \(\mathfrak{h}\) 是 \(\mathfrak{t}\) 的嘉当子代数。我们有 \(\mathfrak{g} = \mathfrak{t} + \mathfrak{r}\) （对 \(\mathfrak{h} \cap \mathfrak{r}\) 在 \(\mathfrak{g}\) 上的伴随表示用 Fitting 分解）。代数 \(\mathfrak{t} \cap \mathfrak{r}\) 是 \(\mathfrak{t}\) 的根基且幂零（利用 §2 习题 5）。

\begin{enumerate}
\def\labelenumi{\alph{enumi})}
\setcounter{enumi}{7}
\tightlist
\item
  假设 k 代数闭。设 h 是 g 的嘉当子代数。存在 g 的 Levi 子代数 s，使得对一切 \(\lambda \in h^{*}\) ，
\end{enumerate}

\[
\mathfrak {g} ^ {\lambda} (\mathfrak {h}) = \left(\mathfrak {g} ^ {\lambda} (\mathfrak {h}) \cap \mathfrak {s}\right) + \left(\mathfrak {g} ^ {\lambda} (\mathfrak {h}) \cap \mathfrak {r}\right).
\]

（沿用 g) 的记号，把 t 的一个满足 \(\mathfrak{h} = (\mathfrak{h} \cap \mathfrak{s}) + (\mathfrak{h} \cap \mathfrak{r})\) 的 Levi 子代数取作 s；由 b) 它存在。）\(^{7}\)

¶ 12) a) 设 g 是可解李代数，G 是 Aut(g) 的有限子群（若 k = R 或 C，相应地取紧子群）。证明存在 g 的在 G 下稳定的嘉当子代数。（对 dim g 作归纳论证，并化归到如下情形：g 是一个幂零代数 g/n 借助一个交换理想 n 的扩张，其中 n 是非平凡单 g/n-模（参见定理 3 的证明）。此时 g 的嘉当子代数构成一个依附于 n 的仿射空间，G 在其上作用。用重心论证得出结论。）

\begin{enumerate}
\def\labelenumi{\alph{enumi})}
\setcounter{enumi}{1}
\tightlist
\item
  设 g 是可解李代数，s 是 \(\text{Der}(\mathfrak{g})\) 的李子代数。假设 s-模 g 半单。证明存在 g 的在 s 下稳定的嘉当子代数。（同一方法。）
\end{enumerate}

¶ 13) 设 g 是李代数，G 是 Aut(g) 的有限子群。假设 G 超可解（《代数学》第一章 §6 习题 26）。证明存在 g 的在 G 下稳定的嘉当子代数。

（对 dim g 作归纳论证。利用习题 12，化归到 g 半单的情形。若 G ≠ \{1\} ，选取 G 的一个素数阶循环正规子群 C（《代数学》，同前引）。由在 C 下不变的元素组成的子代数 s 在 g 中约化（§1 第 5 节），且异于 g。由归纳假设，s 有一个在 G 下稳定的嘉当子代数 a。我们有 s ≠ 0（第一章 §4 习题 21 c)），故 a ≠ 0。a 在 g 中的交换子 z 异于 g 且在 G 下稳定；选取 z 的一个在 G 下稳定的嘉当子代数 h，并证明 h 是 g 的嘉当子代数，参见命题 3 的推论。）

构造 \(\mathfrak{sl}(2,\mathbf{C})\) 的一个有限自同构群，它同构于 \(A_{4}\) （从而可解），且不保持任何嘉当子代数稳定。

\begin{enumerate}
\def\labelenumi{\arabic{enumi})}
\setcounter{enumi}{13}
\item
  *证明超可解有限群 G 的每个不可约复（相应地实）线性表示由 G 的一个子群的 1 次（相应地 1 次或 2 次）表示所诱导。（对李代数 \(\mathfrak{gl}(n,\mathbf{C})\) 与 \(\mathfrak{gl}(n,\mathbf{R}).\) 应用习题 13。）
\item
  假设 k 代数闭。设 g 是李代数，h 是 g 的嘉当子代数，A 是 g 的子集。假设 A 在 g 中稠密（按 Zariski 拓扑）且在 \(\operatorname{Aut}_{e}(\mathfrak{g})\) 下稳定。证明 \(A \cap h\) 在 h 中稠密。（设 X 是 \(A \cap h\) 的闭包，\(U = h - X\) 。假设 \(U \neq \varnothing\) 。用引理 2 的记号，\(U \times \mathfrak{g}^{\lambda_{1}}(\mathfrak{h}) \times \cdots \times \mathfrak{g}^{\lambda_{p}}(\mathfrak{h})\) 在 F 下的像包含 g 的一个非空开子集。由于该像含于 g-A，这与 A 在 g 中稠密矛盾。）
\item
  设 V 是有限维 k-向量空间，g 是 \(\mathfrak{gl}(V)\) 的李子代数。我们将证明下列三个性质等价：(i) g 的嘉当子代数是交换的，且由半单元素组成。
\end{enumerate}

\begin{enumerate}
\def\labelenumi{(\roman{enumi})}
\setcounter{enumi}{1}
\item
  \(\mathfrak{g}\) 的每个正则元都是半单的。
\item
  g 的半单元素在 g 中按 Zariski 拓扑稠密。
\end{enumerate}

\begin{enumerate}
\def\labelenumi{\alph{enumi})}
\item
  证明 (i) \(\Longrightarrow\) (ii) \(\Longrightarrow\) (iii)。
\item
  设 A 是 \(\mathfrak{g}\) 的半单元素之集。证明 A 在 \(\mathrm{Aut}_e(\mathfrak{g})\) 下稳定。
\item
  证明 (iii) \(\Longrightarrow\) (i)。（（附录 I 习题 1）化归到 k 代数闭的情形。若 h 是 g 的嘉当子代数，利用习题 15 证明 \(A \cap h\) 在 h 中稠密；由于 \([x, y] = 0\) 对 \(x \in A \cap h, y \in h\) 成立，推出 h 交换，由此 (i) 立得。）
\item
  假设 k 是 R、C 或特征零的非离散完备超度量域。给 g 赋予由 k 的拓扑所定义的拓扑。证明性质 (i)、(ii)、(iii) 等价于下列性质：
\end{enumerate}

\begin{enumerate}
\def\labelenumi{(\roman{enumi})}
\setcounter{enumi}{3}
\tightlist
\item
  \(\mathfrak{g}\) 的半单元素在 \(\mathfrak{g}\) 中稠密。
\end{enumerate}

（证明 (iv) \(\Longrightarrow\) (iii) 与 (ii) \(\Longrightarrow\) (iv)，参见附录 I 习题 4。）

\begin{enumerate}
\def\labelenumi{\arabic{enumi})}
\setcounter{enumi}{16}
\tightlist
\item
  假设 \(k\) 代数闭。设 \(\mathfrak{g}\) 是李代数，\(\mathfrak{h}\) 是其嘉当子代数，A 是 \(\mathfrak{h}\) 的中心的子集。用 \(\mathrm{E}_{\mathfrak{g}}\) 表示 \(\operatorname{Aut}(\mathfrak{g})\) 的第 2 节中记作 E 的子群。证明：若 \(s\) 是 \(\operatorname{Aut}(\mathfrak{g})\) 的元素且 \(s\mathrm{A} = \mathrm{A}\) ，则存在 \(t \in \mathrm{E}_{\mathfrak{g}}\) 使得 \(t\mathfrak{h} = \mathfrak{h}\) 且 \(t|\mathrm{A} = \mathrm{Id}_{\mathrm{A}}\) ；特别地，\(ts|_{\mathrm{A}} = s|_{\mathrm{A}}\) 。（设 \(\mathfrak{a}\) 是 A 在 \(\mathfrak{g}\) 中的交换子；由于 \(\mathfrak{h}\) 与 \(sh\) 都是 \(\mathfrak{a}\) 的嘉当子代数，存在 \(\theta \in \mathrm{E}_{\mathfrak{a}}\) 使得 \(\theta(s\mathfrak{h}) = \mathfrak{h}\) ；把 \(t\) 取为 \(\mathrm{E}_{\mathfrak{g}}\) 中开拓 \(\theta\) 的一个元素。）
\end{enumerate}

¶ 18) 设 g 是李代数，h 是 g 的嘉当子代数，Ug（相应地 Uh）是 g（相应地 h）的包络代数。Ug 上的线性形式 \(\varphi\) 称为中心的，若它在 {[}Ug, Ug{]} 上为零，即若 \(\varphi(a.b) = \varphi(b.a)\) 对一切 \(a, b \in Ug\) 成立。

\(a)\) 设 \(x, y \in \mathfrak{g}\) 。假设存在 \(s \in \operatorname{Aut}_e(\mathfrak{g})\) 使得 \(s(x) = y\) 。证明

\[
\varphi (x ^ {n}) = \varphi (y ^ {n})
\]

对一切 \(n \in \mathbf{N}\) 与每个中心线性形式 \(\varphi\) （\(\mathrm{Ug}\) 上的）成立。

\begin{enumerate}
\def\labelenumi{\alph{enumi})}
\setcounter{enumi}{1}
\item
  设 \(\varphi\) 是 Ug 上的中心线性形式，它在 Uh 上的限制为零。证明 \(\varphi = 0\) 。（可以假设 k 代数闭。由 a) 推出此时 \(\varphi(x^{n}) = 0\) 对一切 \(n \in N\) 及一切正则的 \(x \in g\) 成立；用稠密性论证去掉正则性假设。）
\item
  证明 \(Ug = [Ug, Ug] + Uh\) 。
\end{enumerate}

\(d\) ) 设 V 是半单 \(\mathfrak{g}\) -模。证明 V 作为 \(\mathfrak{h}\) -模半单。特别地，\(V^{\lambda}(\mathfrak{h}) = V_{\lambda}(\mathfrak{h})\) 对一切 \(\lambda \in \mathfrak{h}^{*}\) 成立。

\begin{enumerate}
\def\labelenumi{\alph{enumi})}
\setcounter{enumi}{4}
\tightlist
\item
  设 \(V'\) 是半单 g-模。假设 V 与 \(V'\) 作为 h-模同构。证明它们作为 g-模同构。（若 \(a \in Uh\) ，注意 \(\operatorname{Tr}(a_{\mathrm{V}}) = \operatorname{Tr}(a_{\mathrm{V}'})\) 。）利用 b) 或 c) 推出 \(\operatorname{Tr}(x_{\mathrm{V}}) = \operatorname{Tr}(x_{\mathrm{V}'})\) 对一切 \(x \in Ug\) 成立，并借助《代数学》第八章得出结论。
\end{enumerate}

若 k 代数闭，则假定“V 与 V' 是 h-同构的”等价于说 \(\dim\mathrm{V}_{\lambda}(\mathfrak{h})=\dim\mathrm{V}_{\lambda}^{\prime}(\mathfrak{h})\) 对一切 \(\lambda\in\mathfrak{h}^{*}\) 成立。

